La predicción bursátil con aprendizaje profundo se ha desarrollado, en su mayor parte,
sobre los mercados en los que cada arquitectura fue diseñada, y casi nunca se comprueba si
la capacidad predictiva se mantiene al trasladar el modelo a un mercado distinto. Este
trabajo estudia esa transferencia entre mercados en un caso concreto. El modelo es
Stockformer, una red de en torno a un millón de parámetros que combina una descomposición
wavelet, atención y un grafo de relaciones entre acciones, calibrada sobre acciones chinas
de clase A. La pregunta es doble: si el modelo conserva su capacidad predictiva al
reentrenarse sobre el S\&P 500 y, en caso de que no la conserve, de dónde sale el retorno
neto que sí se observa en ese mercado.

La primera parte compara Stockformer con modelos más simples (regresiones regularizadas y
árboles potenciados) mediante un mismo marco de evaluación, sobre un mismo panel de datos y
a horizonte diario. La métrica principal es el coeficiente de información de rangos (IC),
que mide la correlación entre el orden de retornos predicho y el realizado. Se revisa
además la posible fuga de información; la única incidencia detectada favorece al modelo
complejo y se reporta por separado. En esta implementación y sobre el universo analizado
(477 componentes del S\&P 500 con cobertura completa de precios), Stockformer no muestra
evidencia robusta de transferibilidad. Su IC fuera de muestra es de $-0{,}003$,
indistinguible de cero, y una regresión Lasso de cinco coeficientes alcanza $+0{,}0238$.
La ventaja de los modelos simples es consistente en signo, pero no alcanza la
significancia estadística al 5\,\%.

La segunda parte estudia el origen del retorno neto de una estrategia semanal neutral al
mercado sobre el mismo universo. Una atribución etapa a etapa muestra que la señal genera
un retorno bruto débil y que la construcción de cartera con control explícito de los
costes de transacción es la etapa que lo conserva en términos netos. Una búsqueda de señal
con criterios fijados de antemano, y validada sobre una muestra de reserva, señala el
momentum residual como el candidato más prometedor. Combinado con esa construcción, el
Sharpe neto pasa de $0{,}38$ a $0{,}72$ ($t=1{,}82$) sobre una ventana de 311 semanas.
Ningún candidato supera el umbral de significancia fijado de antemano, de modo que la
mejora se presenta como evidencia direccional y no como un hallazgo confirmado.

Las conclusiones quedan condicionadas al diseño del estudio: un universo restringido a los
componentes actuales del índice con cobertura completa, datos públicos sin reconstrucción
histórica, una única ejecución del modelo y una función de pérdida adaptada respecto de la
original. Bajo esas condiciones, la contribución es doble. Por un lado, evidencia empírica
de ausencia de transferencia en este caso; por otro, un procedimiento de evaluación y
atribución reproducible que sitúa el origen del retorno neto en una señal simple y en la
construcción de la cartera, no en la complejidad del predictor.